\section{The middle-range field and the denominator on a common family}

% BEGIN A83G1
\paragraph{A83G1. The integral comparison for every inert C and D point.}
Fix the setup and late-family cutoff of I82W1--I82W5 and U83M1.
Take any guarded inert index in row C or D, any good point, any
place above its rational prime, and any of its tested surjections
of crystalline lattices over an allowed coefficient DVR.
The corresponding normalized length of the cokernel on the finite
local cohomology is at most the C/D row entry for this place.

Indeed, if every host position is nonthick, N82I1 gives the
length bound. If every position is thick, T83A6 gives it on the
same surjection. In the mixed case cut the cycle immediately
before each nonthick position. Every arc begins at such a position,
has length one, two or three and has alternating signs by I82W5
and U83M1. In the constant coordinates of S83M1 the complementary
coefficient satisfies the precise pole, residue and filtration
conditions of the scalar lift. There are the following possibilities.
\begin{itemize}
\item S83L2 treats both one-letter signs. For two letters ending
up, S83L2--S83L4 treat the separated and overlapping ranges.
For two letters ending down, S83L3 and D83L2 give the two sides
of the initial-index test.
\item For three letters ending down, put $y=s_{c-2}+S$ and
$b_1=s_{c-1}$. F83S2 applies when $p(y+\eta)\le b_1$,
including equality. When $p(y+\eta)>b_1$, K83J1--K83J3 and
K83V1--K83V2 lift the constant residue and the remaining vector.
\item For three letters ending up, A83U1--A83U3 treat the
two ranges of the initial filtered index, with the same native
band indicator on the cut.
\end{itemize}
All exponents obtained are integers at most $40U+\eta I$, where
$U$ is the local parameter in U83M1 and $I$ is the original band
indicator. A lift at a smaller exponent can be multiplied by
the missing integral power of $p$. Thus every complementary
coordinate lifts at $\lfloor40U+\eta I\rfloor$.
S83M2 applies on each original cut. U83M6 realizes the exceptional
classes in the actual target lattice, and U83M4 sums their
numbers to at most $4dm$ over the cycle. With $E$ the common
exponent there and $e'$ its coarse annihilation exponent, the
whole normalized length is at most $(d+1)rE+4dm e'$. The
coefficient estimates of U83M4 put this quantity inside the
original C/D row with its fixed constants $368,2082,22365$.
This checks the place test at every allowed coefficient extension.
Applying G82I1--G82I2 to these place bounds at the chosen point
gives the original global row inequality there, including the
character and root-difference terms.
% END A83G1

% BEGIN A83G2
\paragraph{A83G2. All fifteen fields hold on the same hub.}
Assume the original $\mathrm{Hypotheses}_2$.
Fix $\mathrm{Fm}$ satisfying $\mathrm{TwGood}$ and
$\mathrm{FactsF}(\mathrm{Fm},C,C_T)$ with $C\ge0$.
Then $\mathrm{ConclF}(\mathrm{Fm})$ holds. Its two constants
are chosen before $x$, the family member, and the local tests.

Choose the free row constant $C_2=122$ before applying F78D2.
Take a common family constant $C_s$ as in G82F4, high enough
for all the fourteen fields constructed there and in its cited
predecessors. Now take the maximum of their cutoffs and those
of the split, inert and arc constructions, for these fixed
constants. Each of the latter is uniform in the point and the
coefficient test by I82W1 and U83M1. At a qualified $x,i$ with
the literal $\mathrm{ShaFinOdd}$ premise, prescribe the member's
exceptional set as in J78H1 and use the resulting F78D2 hub $X$.
It carries all fourteen earlier fields, including the table for
$\mathrm{two}$ in G82F4. The common constant is fixed when that
table and the bound for $X.\mathrm{errorBoundHub}$ are constructed.

For $\mathrm{mid}$, fix a twisted index outside the type primes
which satisfies $\mathrm{RangeAD}$ for $X.\mathrm{exc}$, and a
good point at this index. The range omits the small and ramified
primes as in G82I2. For a split prime apply G82C5. For an inert
prime in row B apply B82R3, for which the chosen $C_2$ suffices.
All other inert rows in this range are C or D, covered by A83G1.
Thus $\mathrm{StarAD}(L,X)$ holds on the same representation,
with its own point, lengths and error. Adjoin this field to
\[
\begin{gathered}
 \mathrm{pos},\ \mathrm{pass},\ \mathrm{mult},\ \mathrm{rank},\
 \mathrm{kill},\ \mathrm{num},\ \mathrm{bb},\\
 \mathrm{exc},\ \mathrm{typ},\ \mathrm{par},\ \mathrm{small},\
 \mathrm{large},\ \mathrm{ram},\ \mathrm{two}.
\end{gathered}
\]
This is $\mathrm{Concl}(\mathrm{Fm}.L(x,i),\mathrm{Fm}.\mathrm{params}
(x,h_x,i),X,C_s)$ for every $x$ beyond the common cutoff and
$h_x:\mathrm{Fm}.x_1<x$, exactly the required family quantifiers.
% END A83G2

% BEGIN A83G3
\paragraph{A83G3. Exporting the full denominator estimate on the given family.}
Under the hypotheses of A83G2 one has $\mathrm{LinkN}(\mathrm{Fm})$.
Write $C_a=C_s$ for its table constant. DNE14 supplies the
export constant $C_{\rm out}=2000(|C_a|+1)^3$. Enlarge the cutoff
to dominate $\mathrm{Fm}.x_1+1$ and the uniform export cutoff at
the given exponent and $C_a$. Fix a later member satisfying
$\mathrm{ShaFinOdd}$. The hub of A83G2 gives its admissible
table by DNE3--DNE4, or equivalently the construction in LN10.
In particular the table uses the actual small-prime witness,
placing its total in the non-place entry at a small prime and
distributing it over the places at $3$ when $3$ is outside the
small set. The large and exceptional terms use this hub's
exceptional set. Apply DNE14 to the table on this member $D$.
For every permitted lift $j$ and under $\mathrm{Summable}
(\mathrm{vvTerm}\,D)$ the conclusion is
\[
 \mathrm{etaSum}(L,j)\le dZ_0(j)\!
 \left[C_{\rm out}\{fm b\,L_{\rm sc}^{\,2}
       +(\log x)^2 N_c\}+\sum_{\ell}\mathrm{vvTerm}(D,\ell)\right],
 \qquad L_{\rm sc}=L.\mathrm{sc}.L .
\]
Here $x=L.\mathrm{sc}.x$ by the late family tie. The left
side sums the original $\eta_J$ over places with a changed
weight, with their logarithmic and ramification weights. The
last summand uses $\mathrm{shaLen}$ of this very $D$ at
$1\le e<t$, with $\log\ell$ at primes and zero off primes.
These identifications are the ones in DNE14, so no rank or
host substitution enters the estimate.

For the selected-design obligation this proves the implication
once the specified family carries the hypotheses of A83G2.
The family in $\mathrm{Des}\,\mathcal C\,3$ already has
$\mathrm{TwGood}$ by its definition. The alternative formal
route applies the analytic estimate to a cut family with
transported character choices, then transports it back. Its
arithmetic inputs include the Galois preconditions and the
fourth tie on the same selected members and the corresponding
good-cut facts. An estimate proved for a family with supplied
$\mathrm{FactsF}$ alone leaves that identification to be shown.
% END A83G3

% BEGIN A83G4
\paragraph{A83G4. A sufficient family-facts statement with the original design quantifiers.}
Under $\mathrm{Hypotheses}_2$ fix the class $\mathcal C$ of the
reconstruction. To obtain its whole denominator contract it suffices to
prove the following statement at $\mathrm{Des}(\mathcal C,3)$:
\[
 \mathrm{Closed}\bigl(\mathrm{Des}(\mathcal C,3),\,
   (\mathrm{Fm},C_H)\mapsto
   \exists C\ge0,\ \exists C_T\in\mathbb N,\
   \mathrm{FactsF}(\mathrm{Fm},C,C_T)\bigr).
\]
In this sufficient condition $C$ and $C_T$ may depend on the
family. Each quantifies over all the members required by
$\mathrm{FactsF}$ on that family.

Indeed, retain the positive $\varepsilon_1$ in this statement.
For each $0<\varepsilon\le\varepsilon_1$ take its very same
transport parameter $K$. Now fix $\mathrm{Fm},C_H$ such that
$\mathrm{Des}(\mathcal C,3)(\mathrm{Fm})$ and
$\mathrm{Fm}.\mathrm{Large}(K,C_H)$. The former gives
$\mathrm{TwGood}(\mathrm{Fm})$ by projection. Specialize the
displayed premise to this pair to obtain $C,C_T$, before applying
A83G2--A83G3 to this fixed family. This yields the constants and
cutoff of $\mathrm{LinkN}(\mathrm{Fm})$ inside the family
quantifiers, with its literal member and summability guards.
Returning $\varepsilon_1$ and $K$ proves the exact
$\mathrm{Closed}(\mathrm{Des}(\mathcal C,3),\mathrm{LinkN})$
conclusion under the displayed premise. This is a direct
sufficient condition on the uncut families. The formal good-cut
route uses $\mathrm{FactsFAlt}$ on a modified family and has its
own local comparison and return map.
% END A83G4
